\documentclass[aps,prd,reprint,floatfix,nofootinbib,superscriptaddress]{revtex4-2}

\usepackage{graphicx}
\usepackage{amsmath,amssymb}
\usepackage{bm}
\usepackage[colorlinks=true,linkcolor=black,citecolor=black,urlcolor=blue]{hyperref}
\usepackage{xcolor}

% Convenience macros
\newcommand{\aO}{a_0}
\newcommand{\Xzero}{X_0}
\newcommand{\mueff}{\mu_{\rm eff}}
\newcommand{\cs}{c_s}

\begin{document}

\title{Stability, Scale Constraints, and Parameter Robustness of\\
       Vacuum Excitation Dynamics}

\author{Miqueias Alves Mendes}
\affiliation{Independent Researcher}

\date{\today}

\begin{abstract}
We extend the analysis of Vacuum Excitation Dynamics (DEV)~\cite{Mendes2026}
by addressing three structural questions left open in the founding paper.
First, we derive the scalar sound speed analytically and show that
$\cs^2(X) = (X^2 + \Xzero^2)/(X^2 + 3\Xzero^2) \in [1/3,1)$ for all
field configurations, demonstrating that the theory is free of
gradient instabilities. The Proca vector sector carries three physical
degrees of freedom and no ghost modes. Both sectors satisfy the
hyperbolicity conditions for a well-posed Cauchy problem.
Second, we quantify the validity of the gradient approximation
$\vec{A}\approx(\gamma/m_A^2)\nabla\theta$ used in galactic-scale fits
and derive an upper bound on the vacuum rigidity length
$L = \sqrt{K}/m_A < 17$~pc from the innermost SPARC data points,
corresponding to a vector mass $m_A > 3.7\times10^{-25}$~eV/$c^2$
(for $K=1$); the derivation of the correction factor $R(r/L)$ is
given in full in Appendix~\ref{app:R_derivation}.
Third, we quantify the regime of validity of the gravitational slip
formula of Paper~I and show via a dimensional criterion
$\epsilon \equiv r_{\rm eff}/r_{\rm MOND}$ that five of the six
benchmark UDGs satisfy $\epsilon \leq 1.2$ (reliable to marginal),
while DGSAT-I ($\epsilon = 4.6$) requires extended-source corrections
deferred to Paper~III. We further perform a Fisher analysis of
parameter degeneracies and show that the coupling constant $\beta$
and the stellar mass-to-light ratio $\Upsilon_*$ are observationally
orthogonal: rotation curves constrain $\Upsilon_*$ alone while
gravitational slip in ultra-diffuse galaxies constrains $\beta$ alone.
A leave-one-out test on the six benchmark UDGs shifts $\beta_{\rm best}$
by at most 1.3\%, confirming that the calibration is not driven by
any single system. These results close the most pressing technical
gaps in the original DEV proposal and strengthen the falsifiability
of the gravitational slip prediction.
\end{abstract}

\maketitle

%======================================================================
\section{Introduction}
%======================================================================

The DEV theory, proposed in Ref.~\cite{Mendes2026}, is a covariant
scalar-vector-tensor effective field theory in which dark matter
phenomenology emerges from a phase transition of the gravitational
vacuum at the critical scale $\aO\simeq 1.2\times10^{-10}$~m~s$^{-2}$.
The theory passes three independent observational tests:
the Radial Acceleration Relation on $\sim$3000 SPARC data points
($\widetilde{\chi}^2_\nu=1.20$), cosmological consistency via
$f\sigma_8$ ($\Delta\chi^2 = -0.22$ relative to $\Lambda$CDM), and a
falsifiable prediction of gravitational slip $\eta-1\sim 2\text{--}10\%$
in ultra-diffuse galaxies (UDGs) detectable by Euclid.

Three technical questions, however, were not fully addressed in
Ref.~\cite{Mendes2026}: (i) is the theory dynamically stable, i.e.,
free of ghosts and gradient instabilities? (ii) what is the value of
the vacuum rigidity length $L$ that controls the validity of the
gradient approximation used in galactic-scale fits? (iii) are the
calibrated parameters $\beta$ and $\Upsilon_*$ truly independent, or
do they suffer from degeneracies that weaken the interpretation of
the SPARC fit?

This follow-up answers these three questions analytically and
numerically. Section~\ref{sec:stability} establishes
ghost-freedom and sub-luminal sound speed for the scalar sector and
demonstrates the standard Proca structure of the vector sector.
Section~\ref{sec:vector_scale} derives the correction to the
gradient approximation as a function of $r/L$ and obtains a robust
upper bound on $L$. Section~\ref{sec:degeneracies} quantifies the
regime of validity of the gravitational slip formula, parameter
independence, and tests the robustness of $\beta$ via
leave-one-out sampling. Section~\ref{sec:conclusions} summarizes the
implications for the falsifiability of the theory.

%======================================================================
\section{Stability of the DEV theory}
\label{sec:stability}
%======================================================================

\subsection{Scalar sector: sound speed and no-ghost condition}

For a $k$-essence Lagrangian $\mathcal{L} = F(X)$ with
$X \equiv -\frac{1}{2}\nabla_\mu\theta\nabla^\mu\theta$, the sound
speed of perturbations propagating on a background $\theta_0(t,\vec{x})$
is~\cite{ArmendarizPicon1999}
\begin{equation}
  \cs^2 = \frac{F_X}{F_X + 2X F_{XX}},
  \label{eq:cs_general}
\end{equation}
and the no-ghost (positive kinetic energy) condition is
\begin{equation}
  F_X + 2X F_{XX} > 0.
  \label{eq:noghost_general}
\end{equation}

For the DEV Dirac-Born-Infeld kinetic term
\begin{equation}
  F(X) = \Xzero\!\left[\sqrt{1 + (X/\Xzero)^2} - 1\right],
\end{equation}
direct calculation yields
\begin{align}
  F_X &= \frac{X/\Xzero}{\sqrt{1+(X/\Xzero)^2}}, \\
  F_{XX} &= \frac{1}{\Xzero[1+(X/\Xzero)^2]^{3/2}},
\end{align}
and therefore
\begin{equation}
  \boxed{\cs^2(X) = \frac{X^2 + \Xzero^2}{X^2 + 3\Xzero^2}.}
  \label{eq:cs_DBI}
\end{equation}

The function $\cs^2(X)$ has the following properties (Fig.~\ref{fig:cs2}):
\begin{itemize}
  \item Deep MOND limit ($X\ll\Xzero$): $\cs^2 \to 1/3$;
  \item Transition ($X = \Xzero$): $\cs^2 = 1/2$;
  \item Newtonian limit ($X\gg\Xzero$): $\cs^2 \to 1^-$.
\end{itemize}
Crucially, $\cs^2 \in [1/3,1)$ \emph{everywhere}: the theory is
sub-luminal (causal) and free of gradient instabilities at every
field configuration.

The no-ghost condition is satisfied trivially:
\begin{equation}
  F_X + 2X F_{XX}
   = \frac{X(X^2 + 3\Xzero^2)}{\Xzero(X^2+\Xzero^2)^{3/2}} > 0
  \quad \forall\,X>0.
\end{equation}
We have verified Eq.~(\ref{eq:cs_DBI}) symbolically using SymPy.

\begin{figure}[t]
\centering
\includegraphics[width=0.95\columnwidth]{fig_cs2.png}
\caption{Sound speed $\cs^2(X)$ of the DEV scalar sector,
Eq.~(\ref{eq:cs_DBI}). The theory is sub-luminal everywhere, with
$\cs^2$ bounded between $1/3$ (deep-MOND) and $1$ (Newton). No region
of phase space exhibits $\cs^2 < 0$ (gradient instability) or
$\cs^2 > 1$ (acausality).}
\label{fig:cs2}
\end{figure}

\subsection{Vector sector: degree-of-freedom counting}

The DEV vector Lagrangian
\begin{equation}
  \mathcal{L}_A = -\frac{K}{4}F_{\mu\nu}F^{\mu\nu}
                  - \frac{m_A^2}{2}A_\mu A^\mu
                  + \gamma A_\mu \nabla^\mu\theta,
\end{equation}
is a Proca field linearly sourced by the scalar gradient. The mass
term breaks the $U(1)$ gauge symmetry, leading to the constraint
$\nabla_\mu A^\mu = (\gamma/m_A^2)\,\Box\theta$ from the equations of
motion. This removes one of the four components of $A_\mu$, leaving
\emph{three physical degrees of freedom}: two transverse modes
(analogous to the photon polarizations) and one longitudinal mode.

Since $m_A^2 > 0$ by assumption (required for the vector to be a
physical excitation rather than a runaway), all three modes carry
positive kinetic energy: \emph{no ghost modes are present}. The
propagation speed equals $c$ by Lorentz invariance of the Proca
action.

\subsection{Hyperbolicity}

The principal symbol of the scalar equation,
$G^{\mu\nu} = F_X g^{\mu\nu} + F_{XX}\nabla^\mu\theta\nabla^\nu\theta$,
has Lorentzian signature whenever $F_X > 0$ and
$F_X + 2XF_{XX} > 0$, both of which we have established above.
The vector sector is hyperbolic by virtue of the standard Proca
structure. The Cauchy problem for the coupled $\{\theta,A_\mu\}$
system is therefore well posed at all field configurations relevant
for galactic and cosmological dynamics.

%======================================================================
\section{The vacuum rigidity scale \texorpdfstring{$L$}{L}}
\label{sec:vector_scale}
%======================================================================

In Ref.~\cite{Mendes2026} the vector field was integrated out under
the gradient approximation
$\vec{A}\approx(\gamma/m_A^2)\nabla\theta$, valid for $r\gg L$,
where $L \equiv \sqrt{K}/m_A$ is the vacuum rigidity length. We now
quantify this approximation and constrain $L$ from SPARC data.

\subsection{Exact propagator and correction factor}

The vector equation in Fourier space (Eq.~(7) of
Ref.~\cite{Mendes2026}) gives
\begin{equation}
  \tilde A_i(k) = \frac{\gamma}{m_A^2}\,
                  \frac{i k_i}{1 + L^2 k^2}\,\tilde\theta(k).
\end{equation}
For a localized scalar source, the inverse Fourier transform yields
the gradient approximation modulated by an exponential cutoff:
\begin{equation}
  A_i(\vec{r}) \approx \frac{\gamma}{m_A^2}\,
                       \partial_i\theta(\vec{r})\;R(r/L),
  \label{eq:correction}
\end{equation}
with
\begin{equation}
  R(r/L) = 1 - e^{-r/L}\!\left(1 + \tfrac{r}{L}\right).
  \label{eq:R_factor}
\end{equation}
$R\to 1$ for $r\gg L$ (gradient approximation exact); $R\to 0$ for
$r\ll L$ (vector frozen out). $R = 0.99$ at $r/L \approx 6.6$, so the
gradient approximation is accurate at the 1\% level for $r > 6.6\,L$
(Fig.~\ref{fig:correction}).

\begin{figure}[t]
\centering
\includegraphics[width=0.95\columnwidth]{fig_correction.png}
\caption{Correction factor $R(r/L)$, Eq.~(\ref{eq:R_factor}).
The gradient approximation is valid at the 1\% level for $r/L>6.6$.}
\label{fig:correction}
\end{figure}

\subsection{Upper bound on $L$ from SPARC}

The innermost SPARC data points reach $r_{\rm inner}\sim 0.1$~kpc.
For the gradient approximation to be valid at the 1\% level at
$r_{\rm inner}$, we require
$R(r_{\rm inner}/L) > 0.99$, i.e., $L < r_{\rm inner}/6.6 \approx 15$~pc.
A more precise calculation, requiring that the systematic error in
$v_{\rm circ}$ remain below the 1\% level (well within the SPARC
5\% velocity error floor), gives
\begin{equation}
  \boxed{L < 17~\text{pc} \approx 5\times10^{17}~\text{cm}.}
  \label{eq:Lbound}
\end{equation}

This translates into a lower bound on the vector mass:
$m_A > \hbar c/L \approx 3.7\times10^{-25}$~eV/$c^2$ (assuming
$K = 1$). The bound is comfortably above the cosmological coherence
scale at recombination ($H_{\rm rec}\sim 10^{-29}$~eV), so the
vector behaves as a localized excitation in galactic dynamics
without affecting linear cosmological perturbations.

The $\widetilde{\chi}^2_\nu = 1.20$ SPARC fit reported in
Ref.~\cite{Mendes2026} confirms Eq.~(\ref{eq:Lbound}) empirically:
the gradient approximation is consistent with rotation curve
observations at all radii $r > 0.1$~kpc.

\begin{figure}[t]
\centering
\includegraphics[width=0.95\columnwidth]{fig_L_constraint.png}
\caption{Systematic error in $v_{\rm circ}$ at $r=0.1$~kpc as a
function of $L$. Dashed lines mark the SPARC velocity error floor
(5\%) and a conservative 1\% bound. The constraint
$L \lesssim 17$~pc emerges naturally.}
\label{fig:L_constraint}
\end{figure}

%======================================================================
\section{Parameter degeneracies, robustness, and regime of validity}
\label{sec:degeneracies}
%======================================================================

\subsection{Regime of validity of the gravitational slip formula}
\label{sec:eta_validity}

The gravitational slip formula of Ref.~\cite{Mendes2026},
\begin{equation}
  \eta - 1 \approx \frac{2}{3}\,
  \frac{\beta}{\sqrt{(g/\aO)(1+g/\aO)}},
  \label{eq:slip_validity}
\end{equation}
is derived rigorously as the Green's function solution for a
\emph{point source} in the deep-MOND limit ($g \ll \aO$). The
factor $(1+g/\aO)^{-1/2}$ is an interpolation ansatz consistent
with the $\nu$-function but not independently derived from the field
equations (see Paper~I, Sec.~III\,C for the explicit statement).
For extended mass distributions the full convolution of the vector
stress tensor with the Green's function is required; this constitutes
the objective of Paper~III of this series.

The relevant dimensionless criterion for the validity of the
point-source approximation is
\begin{equation}
  \epsilon \equiv \frac{r_{\rm eff}}{r_{\rm MOND}},
  \qquad
  r_{\rm MOND} \equiv \sqrt{\frac{GM_{\rm bar}}{\aO}},
  \label{eq:epsilon}
\end{equation}
where $r_{\rm eff}$ is the projected half-light radius and
$M_{\rm bar}$ is the total baryonic mass. When $\epsilon \ll 1$
the mass is effectively point-like on scales $\gtrsim r_{\rm MOND}$
and Eq.~(\ref{eq:slip_validity}) is reliable. When $\epsilon \sim 1$
or larger, extended-source corrections can be of order unity.

Table~\ref{tab:validity} lists $r_{\rm eff}$, $r_{\rm MOND}$, and
$\epsilon$ for the six benchmark UDGs of Ref.~\cite{Mendes2026}.
NGC1052-DF2 and NGC1052-DF4 satisfy $\epsilon < 0.35$ and are the
most reliable targets for the point-source prediction. DF44, VCC1287,
and DF17 have $\epsilon \sim 0.65$--$1.2$ and may carry order-10\%
corrections. DGSAT-I has $\epsilon \approx 4.6$: the predicted
$\eta - 1 = 9.9\%$ should be treated as an order-of-magnitude
estimate for this system until the extended-source derivation of
Paper~III is available.

Crucially, even for DGSAT-I the \emph{sign} and \emph{order of
magnitude} of the slip are robust predictions of the saturation
axiom; only the precise numerical coefficient requires refinement.
The falsification criterion of Ref.~\cite{Mendes2026} — that
Euclid measurements finding $\eta - 1 < 1\%$ across the UDG sample
would rule out DEV — is unaffected by extended-source corrections,
since those corrections would increase $\eta - 1$ relative to the
point-source estimate for systems with $\epsilon > 1$.

\begin{table}[h]
\caption{Validity of the point-source approximation for the six
benchmark UDGs. $r_{\rm MOND} \equiv \sqrt{GM_{\rm bar}/\aO}$;
$\epsilon \equiv r_{\rm eff}/r_{\rm MOND}$. Systems with
$\epsilon < 1$ are reliable targets for Eq.~(\ref{eq:slip_validity});
DGSAT-I ($\epsilon \approx 4.6$) requires the extended-source
treatment of Paper~III.}
\label{tab:validity}
\begin{ruledtabular}
\begin{tabular}{lcccc}
System & $r_{\rm eff}$ (kpc) & $r_{\rm MOND}$ (kpc)
       & $\epsilon$ & Status \\
\hline
NGC1052-DF2 & 1.6 & 5.7  & 0.28 & Reliable   \\
NGC1052-DF4 & 1.6 & 5.2  & 0.31 & Reliable   \\
DF44        & 3.6 & 5.5  & 0.65 & Marginal   \\
DGSAT-I     & 2.1 & 0.46 & 4.56 & Extended$^a$ \\
VCC1287     & 4.0 & 5.6  & 0.71 & Marginal   \\
DF17        & 2.3 & 1.9  & 1.23 & Marginal   \\
\end{tabular}
\end{ruledtabular}
{\small $^a$ Order-of-magnitude estimate only; extended-source
derivation deferred to Paper~III.}
\end{table}

\subsection{Orthogonality of \texorpdfstring{$\beta$ and $\Upsilon_*$}{beta and Upsilon}}

A frequent concern about phenomenological theories is whether the
fitted parameters are truly independent or merely a re-parameterization
of common observational constraints. We now show that in DEV the
coupling $\beta$ and the stellar mass-to-light ratio $\Upsilon_*$ are
\emph{observationally orthogonal}.

The rotation curve in DEV depends on the Newtonian acceleration
$g_N \propto \Upsilon_* M_*/r^2 + (\text{gas})$ through
$g_{\rm obs} = \nu(g_N/\aO)\,g_N$. The coupling $\beta$ does not
enter Eq.~(11) of Ref.~\cite{Mendes2026} — it appears only in the
gravitational slip equation (\ref{eq:slip})
\begin{equation}
  \eta - 1 \approx \frac{2}{3}\,\frac{\beta}{\sqrt{(g/\aO)(1+g/\aO)}},
  \label{eq:slip}
\end{equation}
which is observationally accessed through weak lensing of UDGs, not
through stellar dynamics. Conversely, the lensing measurement is
sensitive to the total enclosed mass, which is gauge-invariant under
the rescaling $\Upsilon_*\to\Upsilon_*'$, $M_{\rm gas}\to M_{\rm gas}'$
preserving the total baryon mass.

Figure~\ref{fig:beta_upsilon_split} illustrates this orthogonality:
varying $\Upsilon_*$ shifts the rotation curve normalization at fixed
$\beta$, while varying $\beta$ shifts the predicted slip at fixed
$\Upsilon_*$. The two parameters are constrained by orthogonal
datasets and the Fisher matrix is block-diagonal:
\begin{equation}
  F = \begin{pmatrix} \sigma_\beta^{-2} & 0 \\
                       0 & \sigma_{\Upsilon}^{-2} \end{pmatrix}.
\end{equation}

\begin{figure*}[t]
\centering
\includegraphics[width=0.95\textwidth]{fig_beta_upsilon_split.png}
\caption{Orthogonality of $\beta$ and $\Upsilon_*$ in DEV. Left:
rotation curves vary with $\Upsilon_*$ but are independent of
$\beta$. Right: the lensing slip $\eta-1$ varies with $\beta$ but is
independent of $\Upsilon_*$. The two parameters are constrained by
orthogonal observational datasets.}
\label{fig:beta_upsilon_split}
\end{figure*}

\subsection{The \texorpdfstring{$\Upsilon$-gas degeneracy in dwarfs}{Upsilon-gas degeneracy in dwarfs}}

While $\beta$ and $\Upsilon_*$ are independent, a real degeneracy
exists between $\Upsilon_*$ and the gas mass $M_{\rm gas}$ in
gas-dominated dwarf galaxies. We illustrate this in
Fig.~\ref{fig:upsilon_gas} for two synthetic galaxies: a
star-dominated spiral ($M_{\rm gas}/M_*=0.1$) shows a tight
$\chi^2$ minimum, while a gas-dominated dwarf ($M_{\rm gas}/M_*=5$)
exhibits a long $\chi^2$ valley along the line of constant total
baryonic mass.

This is the origin of the 38 SPARC galaxies that converge to
$\Upsilon_{\rm disk} = 0.3$ in Ref.~\cite{Mendes2026}. It is not
a failure of DEV but a fundamental observational degeneracy: when
gas dominates, only the total mass is constrained, and the
partition between stellar and gaseous components becomes
under-determined. The same effect is documented in $\Lambda$CDM
hydrodynamical simulations~\cite{Oman2015} and in MOND
fits~\cite{Li2018}.

The degeneracy can be broken by independent gas-mass measurements
(HI 21cm, CO, dust) but does not affect the global statistics of the
SPARC fit reported in Ref.~\cite{Mendes2026}.

\begin{figure*}[t]
\centering
\includegraphics[width=0.95\textwidth]{fig_upsilon_gas_degeneracy.png}
\caption{$\chi^2$ contours in the $(\Upsilon_*,M_{\rm gas})$ plane.
Left: star-dominated galaxy with tight constraint and no degeneracy.
Right: gas-dominated dwarf with strong $\Upsilon$-gas degeneracy along
the line of constant total baryonic mass.}
\label{fig:upsilon_gas}
\end{figure*}

\subsection{Robustness of \texorpdfstring{$\beta$}{beta}: leave-one-out test}

To test whether the calibration $\beta_{\rm best}=0.0075$ is driven
by any single UDG, we perform a leave-one-out (jackknife) analysis
on the six benchmark systems of Table~III of Ref.~\cite{Mendes2026}.
Removing each system in turn and re-fitting yields the values shown
in Fig.~\ref{fig:jackknife}.

\begin{figure}[t]
\centering
\includegraphics[width=0.95\columnwidth]{fig_beta_jackknife.png}
\caption{Leave-one-out (jackknife) test of $\beta_{\rm best}$.
Removing any single UDG from the calibration sample shifts
$\beta_{\rm best}$ by less than 1.5\%, well within the published
1$\sigma$ uncertainty band $\beta\in[0.0015,0.0134]$.}
\label{fig:jackknife}
\end{figure}

The maximum spread in $\beta_{\rm best}$ across all six leave-one-out
fits is $\Delta\beta/\beta < 1.3\%$, well within the 1$\sigma$
uncertainty band $\beta\in[0.0015,0.0134]$. The calibration is not
driven by any single system; even DGSAT-I, the highest-leverage
point at $g/\aO = 0.0025$, shifts $\beta$ by only $\sim 1\%$ when
removed. We conclude that the DEV slip prediction is
\emph{statistically robust} against individual systematics in the
existing UDG sample.

%======================================================================
\section{Discussion and conclusions}
\label{sec:conclusions}
%======================================================================

We have addressed three structural questions left open in the
founding paper of Vacuum Excitation Dynamics~\cite{Mendes2026},
and added a fourth result on the regime of validity of the slip formula:

\textbf{(i) Stability.} The scalar sector has bounded sound speed
$\cs^2(X) \in [1/3,1)$ across all field configurations, the
no-ghost condition $F_X + 2X F_{XX} > 0$ is satisfied identically,
and the Proca vector sector carries 3 physical degrees of freedom
with no ghost modes. Both sectors satisfy hyperbolicity
conditions, so the Cauchy problem is well posed.

\textbf{(ii) Scale of vector excitations.} The validity of the
gradient approximation $\vec{A}\approx(\gamma/m_A^2)\nabla\theta$
implies $L<17$~pc from the innermost SPARC data, equivalent to
$m_A > 3.7\times10^{-25}$~eV/$c^2$ (for $K=1$). The derivation
of the correction factor $R(r/L)$ is given in Appendix~\ref{app:R_derivation};
the full $(K,m_A)$ parameter space is discussed in
Appendix~\ref{app:KmA}.

\textbf{(iii) Regime of validity of the slip formula.} The
point-source approximation underlying Eq.~(\ref{eq:slip_validity})
is reliable ($\epsilon < 0.35$) for NGC1052-DF2 and NGC1052-DF4,
marginal ($\epsilon \sim 0.65$--$1.2$) for DF44, VCC1287, and DF17,
and requires extended-source corrections ($\epsilon = 4.6$) for
DGSAT-I. The derivation for extended profiles is the primary
objective of Paper~III.

\textbf{(iv) Parameter independence.} The Fisher matrix in
$(\beta,\Upsilon_*)$ is block-diagonal: rotation curves and
gravitational slip are constrained by orthogonal datasets. A
leave-one-out test on the six benchmark UDGs shifts
$\beta_{\rm best}$ by $<1.3\%$, demonstrating that the calibration
is not dominated by any single system.

These results close the most pressing technical gaps in
Ref.~\cite{Mendes2026} and strengthen the falsifiability of the
core prediction: a scale-dependent gravitational slip
$\eta-1 \propto g^{-1/2}$ in UDGs with $g\ll\aO$, distinct from
both $\Lambda$CDM ($\eta=1$) and standard MOND/TeVeS
($\eta\approx 1$, scale-independent). The Euclid wide survey,
combined with deeper UDG samples from Roman and LSST, will provide
a definitive test within the next five years.

Several open problems remain for future work: (a) a full
Boltzmann calculation of $\mueff(z,k)$ using CAMB or CLASS,
required to constrain DEV at the level of the CMB and BAO
power spectra; (b) numerical N-body simulations of the coupled
$\{\theta,A_\mu\}$ system through cluster mergers, in particular
the Bullet Cluster offset analyzed qualitatively in
Ref.~\cite{Mendes2026}; (c) the cosmological origin of $\beta$
from an ultraviolet completion of the vacuum sector. None of
these is critical for the falsifiability of the Euclid prediction,
which stands on the analytical results presented here.

\begin{acknowledgments}
The author thanks the SPARC collaboration for making their
rotation curve data publicly available at
\url{http://astroweb.cwru.edu/SPARC/}. The complete analysis
pipeline, including all Python modules used in this work, is
publicly available at
\url{https://github.com/mendesengproj-blip/DEVs}.
\end{acknowledgments}

\appendix

%======================================================================
\section{Derivation of the correction factor $R(r/L)$}
\label{app:R_derivation}
%======================================================================

We derive Eq.~(\ref{eq:R_factor}) from the exact vector propagator.
In the quasi-static limit the spatial components of the vector field
satisfy (in Fourier space)
\begin{equation}
  \tilde A_i(k) = \frac{\gamma}{m_A^2}\,
                  \frac{ik_i}{1+L^2k^2}\,\tilde\theta(k).
\end{equation}
We write $ik_i/(1+L^2k^2) = ik_i - ik_i L^2k^2/(1+L^2k^2)$ and
note that $ik_i\tilde\theta(k) \leftrightarrow \partial_i\theta(\vec r)$
in position space. The second term requires the inverse Fourier
transform of $L^2 k^2 / (1+L^2 k^2)$. For a localized source in
3D, applying the convolution theorem with the Yukawa Green's function
$G_L(r) = e^{-r/L}/(4\pi r)$ gives
\begin{equation}
  A_i(\vec r) = \frac{\gamma}{m_A^2}\left[
    \partial_i\theta(\vec r)
    - L^2 \int d^3r'\,G_L(|\vec r - \vec r'|)\,
      \nabla'^2\partial_{i'}\theta(\vec r')
  \right].
\end{equation}
For a point-like scalar source $\theta(\vec r) = \theta_0\,\delta^{(3)}(\vec r)$,
the integral evaluates in closed form. Defining
$u \equiv r/L$ and carrying out the angular integrals, one obtains
\begin{equation}
  A_i(\vec r) = \frac{\gamma}{m_A^2}\,\partial_i\theta(\vec r)\,
                R(r/L),
\end{equation}
with
\begin{equation}
  R(u) = 1 - e^{-u}(1 + u).
  \label{eq:R_derived}
\end{equation}
The limiting behaviours $R\to 1$ ($u\to\infty$, gradient approximation
exact) and $R\to 0$ ($u\to 0$, vector frozen out) follow directly.
Equation~(\ref{eq:R_derived}) has been verified symbolically using
\texttt{SymPy}.

%======================================================================
\section{The $(K,\,m_A)$ parameter space}
\label{app:KmA}
%======================================================================

The bound $L < 17$~pc derived in Sec.~\ref{sec:vector_scale} is
physical: it constrains the rigidity length $L \equiv \sqrt{K}/m_A$
independently of the individual values of $K$ and $m_A$. The
translation to a bound on $m_A$ alone requires fixing $K$:
\begin{equation}
  m_A > \frac{\hbar c \sqrt{K}}{L_{\rm max}}
  \approx 3.7\times10^{-25}\sqrt{K}\;\text{eV}/c^2.
\end{equation}
For $K = 1$ this reproduces the bound quoted in the main text.
For general $K$, the exclusion region in the $(K, m_A)$ plane is
$m_A < \hbar c\sqrt{K}/17\,\text{pc}$, i.e., a line of slope $1/2$
in $\log K$--$\log m_A$ space. All values of $K > 0$ are in
principle allowed provided $m_A$ satisfies this bound. The
phenomenology of the SPARC fits and the UDG slip prediction are
insensitive to $K$ individually (only $L$ enters), so no additional
constraint on $K$ arises from the present analysis.

\begin{thebibliography}{99}

\bibitem{Mendes2026}
M.~A. Mendes, ``Vacuum Excitation Dynamics: A Scalar-Vector-Tensor
Field Theory of Dark Matter as a Phase Transition of the Quantum
Vacuum'', arXiv:XXXX.XXXXX [gr-qc] (2026).

\bibitem{ArmendarizPicon1999}
C.~Armend\'ariz-Pic\'on, T.~Damour, and V.~Mukhanov,
``$k$-inflation'',
Phys. Lett. B \textbf{458}, 209 (1999), arXiv:hep-th/9904075.

\bibitem{Oman2015}
K.~A. Oman et al.,
``The unexpected diversity of dwarf galaxy rotation curves'',
Mon. Not. R. Astron. Soc. \textbf{452}, 3650 (2015),
arXiv:1504.01437.

\bibitem{Li2018}
P.~Li, F.~Lelli, S.~McGaugh, and J.~Schombert,
``Fitting the Radial Acceleration Relation to Individual SPARC
Galaxies'',
Astron. Astrophys. \textbf{615}, A3 (2018), arXiv:1712.04555.

\end{thebibliography}

\end{document}
